\documentclass[11pt]{article}
\usepackage[T1]{fontenc}
\usepackage{lmodern}
\usepackage{amsmath,amssymb,amsthm}
\usepackage[margin=1in]{geometry}
\usepackage{microtype}
\usepackage{hyperref}
\hypersetup{colorlinks=true,linkcolor=blue,citecolor=blue,urlcolor=blue,
  pdftitle={Odd Euler characteristic for flag-no-square four-manifolds},
  pdfauthor={Kaiwen Jin}}
\newtheorem{theorem}{Theorem}
\newtheorem{proposition}[theorem]{Proposition}
\newtheorem{lemma}[theorem]{Lemma}
\newcommand{\lk}{\operatorname{lk}}
\newcommand{\st}{\operatorname{st}}
\newcommand{\dist}{d_{\mathrm G}}
\title{Odd Euler characteristic for flag-no-square four-manifolds}
\author{Kaiwen Jin\\\small\texttt{yui.ui.siki@gmail.com}}
\date{October 14, 2026}
\begin{document}
\maketitle

\begin{abstract}
We construct a closed connected PL four-manifold with odd Euler
characteristic that admits a finite flag-no-square triangulation.
Here flag means that every clique is a simplex, and no-square excludes
induced four-cycles in the graph. The construction identifies antipodal
points on the boundary of a deleted vertex star. In a regular hyperbolic
tessellation, each graph edge outside that vertex changes the radial
direction by at most $\pi/5$; this prevents the identification from creating
a short nonclosing lift. A twisted double makes both the manifold
condition and the global flag-no-square property explicit.
\end{abstract}

\section{The existence question}

A simplicial complex is \emph{flag} if every finite set of pairwise adjacent
vertices spans a simplex. It is \emph{flag-no-square} if, in addition,
its graph contains no induced cycle of length four. Such triangulations
connect manifold topology with hyperbolic right-angled Coxeter groups;
see Przytycki and \'{S}wi\k{a}tkowski~\cite{PS}.
We consider closed connected four-manifolds, where closed means compact
without boundary and orientability is not required.

Przytycki and \'{S}wi\k{a}tkowski~\cite[Remark~4.4(2)]{PS} give
flag-no-square four-manifolds as finite quotients of a regular hyperbolic
tessellation. Kalmanovich, Nevo and Sorcar~\cite{KNS} introduce
star-connected sums and star handles and obtain every sufficiently large
even Euler characteristic. Their Problem~4.1 asks whether odd Euler
characteristic can occur. The operations above change Euler characteristic
by an even integer. We use a boundary identification that changes it by one.

\begin{theorem}\label{thm:main}
There exists a closed connected PL four-manifold with odd Euler
characteristic admitting a finite flag-no-square triangulation.
\end{theorem}

The difficult point in identifying a boundary with itself is the possible
creation of short cycles through the exterior of the removed ball.
We first give a folding criterion in terms of distance in the entire
antistar. A radial estimate in the regular hyperbolic tessellation supplies
this distance bound. A distant star-connected sum of two copies supplies
an even starting Euler characteristic, so no parity theorem for hyperbolic
manifolds is needed.

For a complex $N$, we write $\dist$ for distance in its graph,
$\lk_v N$ for the link of a vertex $v$, and $N-v$ for the induced
subcomplex on the remaining vertices. Euler characteristic is denoted
by $\chi$ and computed by alternating face counts.

\section{Folding a vertex boundary}

We will realize a folded antistar as a quotient of a closed manifold.
The following observation explains the required distance threshold.

\begin{lemma}\label{lem:quotient}
Let $E$ be a flag-no-square simplicial complex with a simplicial
involution $t$. Suppose
\[
  \dist(x,t x)\geq 5
  \qquad\text{for every vertex }x\in E.
\]
The simplicial quotient $Q=E/\langle t\rangle$ is flag-no-square.
The quotient map preserves vertex stars and identifies $|Q|$
with the topological quotient $|E|/\langle t\rangle$.
If the involution is free, this map is a two-sheeted covering.
\end{lemma}

\begin{proof}
Define faces of $Q$ to be images of faces of $E$.
Two vertices in the star of $x$ cannot have the same image:
their distance is at most two, which contradicts the displacement
assumption if they are interchanged by $t$.
There are also no additional edges between the images of neighbors of
$x$. Indeed, if $u,z$ are such neighbors and an edge between their
images lifts as $u\,t z$ rather than $uz$, then
$z,x,u,t z$ is a walk of length three, again a contradiction.

Every simplex containing the image of $x$ has a unique lift containing
$x$, by aligning its chosen lift of that vertex. Thus the quotient
is simplicially isomorphic on each full vertex star. These local
isomorphisms also show that its realization is the topological quotient
and that a free action gives a covering.

Lift a clique of $Q$ by fixing one of its vertices and lifting all
the others into its star. The adjacency check above makes this a clique
of $E$, so it is a simplex.
A four-cycle in $Q$ lifts to a walk that either closes or joins
$x$ to $t x$. The latter is excluded by the distance bound.
An induced four-cycle that closes would be an induced four-cycle
in $E$, since an upstairs diagonal would project to a downstairs
diagonal. This proves the assertion.
\end{proof}

\begin{proposition}\label{prop:fold}
Let $N$ be a finite flag-no-square triangulation of a closed connected
orientable PL four-manifold. Let $v$ be a vertex whose link
$B=\lk_v N$ is the boundary of a regular $600$-cell, and let
$\tau:B\to B$ be its antipodal map. Put $A=N-v$.
Suppose
\begin{equation}\label{eq:antistar}
  d_{A}(b,\tau b)\geq 5
  \qquad\text{for every vertex }b\in B.
\end{equation}
Then identifying $b$ with $\tau b$ on the boundary of $A$
gives a finite flag-no-square triangulation of a closed connected
PL four-manifold $Q$, with
\[
  \chi(Q)=\chi(N)-1.
\]
\end{proposition}

\begin{proof}
The antipodal map of the $600$-cell is a simplicial involution
acting freely on its three-sphere. For example, its central symmetry is
visible in the usual coordinates~\cite[Example~1.27]{Heckman}.
Take two copies $A_1,A_2$ and form
\[
  E=A_1\cup_\tau A_2,
  \qquad b\in B_1\ \sim\ \tau b\in B_2.
\]
This is the star-connected sum of two copies of $N$. The antipodal
map has degree $+1$ on $S^3$; choosing opposite orientations on the
two copies makes the boundary gluing orientation reversing.
Thus $E$ is a closed connected PL four-manifold by
\cite[Lemma~1.3]{KNS}. We recall the combinatorial reason that it
remains flag-no-square, since it also controls short paths below.

If $y\in A\setminus B$, its neighbors in $B$ form a clique.
Otherwise two nonadjacent such neighbors $b,c$ would give the
induced square $b,v,c,y$ in $N$.
The complexes $A$ and $B$ are induced subcomplexes of $N$.
A clique in $E$ lies in one copy, because there are no edges
between the two exteriors, so $E$ is flag.
An induced square using both exteriors would have two opposite
exterior vertices and two nonadjacent boundary vertices. One of the
exterior vertices would then violate the clique observation.
Hence $E$ is no-square.

Swapping the copies with their original labels defines an involution
$t:E\to E$. It is well defined because $\tau^2=1$.
On the common boundary, using first-copy labels, it acts as $\tau$.
It is free, since it interchanges the exteriors and acts antipodally
on the common three-sphere.

We claim $\dist(x,t x)\geq 5$ for every vertex $x$ of $E$.
Suppose a path $P$ of at most four edges joins $x$ to $t x$.
It meets the common boundary $B$. Choose a boundary vertex $b$
on it. Follow the suffix of $P$ from $b$ to $t x$, then the
$t$-image of the prefix from $x$ to $b$. This is a walk of the
same length from $b$ to $t b=\tau b$.

Decompose this walk into boundary edges and excursions with endpoints
on $B$ and all internal vertices in one exterior. Each excursion
has at least two edges, so there are at most two.
With no excursion, the walk lies in $B$ and contradicts
\eqref{eq:antistar}.
With one excursion, read every boundary edge using the labels of
the copy containing that excursion. This is legitimate because
$\tau$ preserves all boundary edges. The resulting walk in $A$
joins $b$ to $\tau b$, or $\tau b$ to $b$, and again contradicts
\eqref{eq:antistar}.
With two excursions, both have length two. The clique observation
replaces each by a boundary edge, giving a boundary walk of length
at most two from $b$ to $\tau b$. This is impossible as well.

Lemma~\ref{lem:quotient} now applies. Its covering assertion shows
that $Q$ is an actual closed connected PL four-manifold. Its
description as $E/\langle t\rangle$ is exactly the identification
of antipodal boundary points of $A$.

The boundary $B$ has face numbers $(120,720,1200,600)$,
so $\chi(B)=0$. The closed vertex star is a cone of Euler
characteristic one. Face-count inclusion-exclusion gives
$\chi(A)=\chi(N)-1$ and $\chi(E)=2\chi(A)$.
Every face of $E$ has an orbit of size two: a face fixed by $t$
would contain a vertex and its image at distance at most one.
By the star isomorphisms, faces of $Q$ are precisely these orbits.
Therefore $\chi(Q)=\chi(E)/2=\chi(N)-1$.
\end{proof}

\section{Antipodal paths in the regular tessellation}

Let $T$ be the regular simplicial tessellation of $\mathbb H^4$
with Schlaefli symbol $\{3,3,3,5\}$. It is flag-no-square, its
vertex links are regular $600$-cells, and its reflection group is
$W=[3,3,3,5]$, the reversal of $[5,3,3,3]$; see
\cite[Remark~4.4(2)]{PS} and \cite[p.~3]{KNS}.
Write $e>0$ for the common hyperbolic length of its edges.

\begin{lemma}\label{lem:separation}
Distinct vertices of $T$ have hyperbolic distance at least $e$.
\end{lemma}

\begin{proof}
Take a fundamental flag orthoscheme $C$ of the reflection tessellation,
with a honeycomb vertex $p$ at the intersection of mirrors numbered
$1,2,3,4$. Its stabilizer
$H=\langle r_1,r_2,r_3,r_4\rangle$ is the finite spherical
reflection group $H_4$. Reflection $r_0$ in the opposite wall
sends $p$ to an adjacent honeycomb vertex, so that wall has
distance $e/2$ from $p$.

Group the chambers incident to $p$ into
$D_p=\bigcup_{h\in H}hC$.
The chamber stabilizer and fundamental-domain statements used here
are \cite[Theorem~4.17, Corollary~4.19 and Proposition~5.8]{Heckman}.
The point $p$ lies in the interior of $D_p$, since its incident
chambers fill the tangent sphere.
Every boundary face of $D_p$ lies on a mirror $h\operatorname{Fix}(r_0)$.
Faces on the other four mirrors are paired within $H$; more generally,
all chambers incident to a face not involving the opposite wall
belong to the same finite stabilizer coset.
Each remaining mirror has distance $e/2$ from $p$.
The open ball of radius $e/2$ about $p$ is consequently contained
in $D_p$: a path leaving this region has to cross its boundary.

The regions $wD_p$, indexed by cosets $wH$ and hence by honeycomb
vertices $wp$, have disjoint interiors. An open overlap would contain
a chamber interior from both cosets, forcing the cosets to agree.
Each region contains the open ball of radius $e/2$ about its vertex.
These balls have disjoint interiors, proving the claimed separation.
\end{proof}

\begin{lemma}\label{lem:paths}
Fix a vertex $p$ of $T$ and let $\tau$ act antipodally on its
spherical $600$-cell link. Every graph path from a link vertex
$b$ to $\tau b$ avoiding $p$ has at least five edges.
\end{lemma}

\begin{proof}
An equilateral triangular face of $T$ has angle $\alpha=\pi/5$
at each vertex. Here is a direct computation of this angle.
Five regular four-simplices meet at each triangular face,
so their dihedral angle is $2\pi/5$.
Set $c=\cos(2\pi/5)$ and $\phi=(1+\sqrt5)/2$.
The normal Gram matrix of one such simplex is
$(1+c)I-cJ$, with five rows, where $J$ has all entries equal
to one. Its inverse has diagonal
$(1-3c)/((1+c)(1-4c))$ and off-diagonal
$c/((1+c)(1-4c))$. The normalized dual timelike vertex
vectors therefore give
\[
  \cosh e=\frac{c}{1-3c}=2\phi+1,
  \qquad
  \cos\alpha=\frac{\cosh e}{\cosh e+1}
             =\frac{\phi}{2}=\cos\frac{\pi}{5}.
\]

For an edge $xy$ not incident to $p$, put
$r=d_{\mathbb H}(p,x)$ and $s=d_{\mathbb H}(p,y)$.
Both are at least $e$ by Lemma~\ref{lem:separation}.
Let $\theta$ be the angle between the radial directions to $x,y$.
In the Lorentz hyperboloid model~\cite[Section~5.1]{Heckman},
write
$x=\cosh(r)p+\sinh(r)u$ and
$y=\cosh(s)p+\sinh(s)z$ with unit tangent vectors $u,z$,
so $\langle u,z\rangle=\cos\theta$.
Since $-\langle x,y\rangle=\cosh e$, we obtain
\begin{align*}
 \sin^2\frac{\theta}{2}
 &=\frac{\cosh e-\cosh(r-s)}{2\sinh r\,\sinh s}\\
 &\leq\frac{\cosh e-1}{2\sinh^2 e}
  =\frac{1}{2(\cosh e+1)}
  =\sin^2\frac{\alpha}{2}.
\end{align*}
Thus each graph edge outside $p$ changes the direction in the
unit tangent three-sphere by at most $\pi/5$.
The directions to $b,\tau b$ are antipodal, at spherical distance
$\pi$. The spherical triangle inequality along a path of $k$
edges gives $\pi\leq k\pi/5$, hence $k\geq 5$.
\end{proof}

\section{A finite seed and the odd-characteristic quotient}

The tessellation $T$ has orientation-preserving finite-index
torsion-free subgroups $K<W$ with arbitrarily large graph
displacement:
\[
 \inf_{\substack{a\in V(T)\\1\ne k\in K}}\dist(a,ka)>R
 \quad\text{for any prescribed integer }R.
\]
Their quotients $Y=T/K$ are finite closed connected orientable
PL four-manifolds with flag-no-square triangulations. This is the
construction of \cite[Remark~4.4(2)]{PS}, with the graph-displacement
form and orientation choice explained in \cite[p.~3]{KNS}.

\begin{proof}[Proof of Theorem~\ref{thm:main}]
Choose $K$ with nonidentity displacement at least $31$.
Choose $p,c\in V(T)$ with $\dist(p,c)=10$; this is possible
because $T$ is connected, locally finite and infinite.
Their images $v,w$ in $Y$ have graph distance ten as well.
Otherwise a shorter quotient path, compared with the original
length-ten path, would give a nonidentity displacement of $p$
by at most nineteen.

Perform the star-connected sum of two copies of $Y$ at their
copies of $w$, matching the $600$-cell links by an
orientation-reversing reflection. Denote the resulting finite
flag-no-square triangulation by $N$, and retain $v$ in the first copy.
By \cite[Lemma~1.3]{KNS}, it is a closed connected orientable
PL four-manifold and
\[
  \chi(N)=2\chi(Y)-2.
\]
The link $B=\lk_v N$ is unchanged from $Y$ and has its antipodal
map $\tau$.

We verify \eqref{eq:antistar} for $A=N-v$.
The connected-sum boundary $\lk_w Y$ is at distance at least
nine from $v$ in the original copy of $Y$.
A path starting at a neighbor $b$ of $v$ needs at least eight
edges to encounter that boundary for the first time.
Consequently a path of at most four edges from $b$ to $\tau b$
in $A$ would lie entirely in the first copy of $Y$.

Lift this path to $T$, starting at the chosen lift of $b$
adjacent to $p$. Adding the two incident edges at $v$ gives
a loop of length at most six in $Y$.
Its lift must return to $p$, since every nonidentity element
of $K$ displaces $p$ by at least $31$.
The endpoint of the lifted middle path is therefore the actual
antipode of its initial link vertex. The lift avoids $p$
because the quotient path avoids $v$.
This contradicts Lemma~\ref{lem:paths}.
Thus \eqref{eq:antistar} holds for every $b\in B$.

Apply Proposition~\ref{prop:fold}. The resulting finite
flag-no-square complex $Q$ triangulates a closed connected
PL four-manifold and satisfies
\[
  \chi(Q)=\chi(N)-1=2\chi(Y)-3,
\]
which is odd. This proves the theorem.
\end{proof}

\paragraph{AI assistance.}
OpenAI Codex assisted with developing the construction, checking the
argument and preparing this text. The result is submitted for
independent mathematical review.

\begin{thebibliography}{9}
\bibitem{PS}
P.~Przytycki and J.~\'{S}wi\k{a}tkowski,
\emph{Flag-no-square triangulations and Gromov boundaries in dimension~3},
Groups, Geometry, and Dynamics \textbf{3} (2009), 453--468.
Revised author manuscript:
\url{https://www.mimuw.edu.pl/~pprzytyc/3tree_revisedb.pdf}.

\bibitem{KNS}
D.~Kalmanovich, E.~Nevo and G.~Sorcar,
\emph{On flag-no-square $4$-manifolds},
arXiv:2310.13495v1 (2023),
\url{https://arxiv.org/abs/2310.13495v1}.
The lemma and page numbers cited here refer to this version.
Published as \emph{On Flag-No-Square 4-Manifolds},
Discrete \& Computational Geometry \textbf{76} (2026), 298--312,
doi:\href{https://doi.org/10.1007/s00454-025-00774-x}
{10.1007/s00454-025-00774-x}.

\bibitem{Heckman}
G.~Heckman,
\emph{Coxeter Groups}, unfinished lecture notes,
October~10, 2018,
\url{https://www.math.ru.nl/~heckman/CoxeterGroups.pdf}.
\end{thebibliography}
\end{document}
